%%
%% fall-lecture05.tex
%% 
%% Made by Alex Nelson <pqnelson@gmail.com>
%% Login   <alex@lisp>
%% 
%% Started on  2026-10-08T08:43:28-0700
%% Last update 2026-10-08T08:43:28-0700
%% 

\lecture[]{}

\subsection{``Einstein'' summation convention}

\begin{remark}
Just an aside to the reader (probably ``future me''): Engineers use
\textbf{Euclidean summation convention} (where we sum over \emph{any} repeated
indices, without caring if they ``live upstairs'' [are superscripts]
or ``live downstairs'' [are subscripts]) but insist this is ``Einstein
summation convention'' (which requires one index to live upstairs and
the other index to live downstairs). This is unfortunate.
\end{remark}

\begin{node}
Vectors $\vec{a}\in\RR^{3}$ may be written as
\begin{equation}
\vec{a} = \begin{bmatrix}a_{1}\\
a_{2}\\
a_{3}
\end{bmatrix}
\end{equation}
Similarly for $\vec{b}$ we may write it out as a column vector. Then
we have
\begin{equation}
\vec{a}\cdot\vec{b}=\sum^{3}_{i=1}a_{i}b_{i}
\end{equation}
and the convention adopted is to simply omit the sum
\begin{equation}
\vec{a}\cdot\vec{b}=\sum^{3}_{i=1}a_{i}b_{i}=a_{i}b_{i}.
\end{equation}
When two indices are repeated, they are summed over. The choice of
index is irrelevant, it becomes a ``dummy variable'' (like the
variable of integration).
\end{node}

\begin{example}
The gradient of $f$
\begin{subequations}
  \begin{align}
\vec{\nabla}f &= \partial_{x}f\,\vec{e}_{x} +
\partial_{y}f\,\vec{e}_{y} + \partial_{z}f\,\vec{e}_{z}\\
&= \partial_{i}f\,\vec{e}_{i}.
  \end{align}
\end{subequations}
\end{example}

\begin{example}
The Laplacian of a scalar function $f$ is
\begin{equation}
\nabla^{2}f = \partial_{i}\partial_{i}f.
\end{equation}
\end{example}

\begin{example}
The divergence of a vector
\begin{equation}
\vec{\nabla}\cdot\vec{f}=\partial_{i}f_{i}.
\end{equation}
Observe that $\partial_{j}f_{i}$ is a $3\times3$ matrix.
\end{example}

\begin{example}
The Kronecker delta
\begin{equation}
\delta_{i,j} = \begin{cases}1 & \mbox{if }i=j\\
0 & \mbox{if }i\neq j
\end{cases}
\end{equation}
\end{example}

\begin{example}
Matrix multiplication $\vec{b}=M\vec{a}$ in index notation
\begin{equation}
b_{i}=M_{ij}a_{j}.
\end{equation}
\end{example}

\begin{example}
Coordinate transformations from $\widehat{x}_{i}$ to $\widehat{x}'_{i}$
is handled by a rotation matrix $R_{ij}$ giving us
\begin{equation}
\vec{x}'=R\vec{x}
\end{equation}
or
\begin{equation}
x'_{i}=R_{ij}x_{j}
\end{equation}
This has the property
\begin{equation}
RR^{\mathsf{T}}=I
\end{equation}
which in index notation
\begin{equation}
R_{ik}(R^{\mathsf{T}})_{kj}=R_{ik}R_{jk}=\delta_{i,j}.
\end{equation}
\end{example}

\begin{example}
The Levi--Civita symbol
\begin{equation}
\varepsilon_{ijk} = \begin{cases}
  0 & \mbox{if $i$, $j$, $k$ repeat}\\
  1 & \mbox{if $(i,j,k)$ is an even permutation}\\
  -1 & \mbox{if $(i,j,k)$ is an odd permutation}
\end{cases}
\end{equation}
Then the cross product
\begin{equation}
(\vec{a}\times\vec{b})_{k}=\varepsilon_{ijk}a_{i}b_{j}.
\end{equation}
\end{example}

\begin{example}
The divergence theorem in vector notation
\begin{equation}
\oint_{A}\vec{f}\cdot\UnitNormalVec{n}\,\D A =
\iiint_{V}\vec{\nabla}\cdot\vec{f}\,\D V
\end{equation}
becomes, in index notation,
\begin{equation}
\oint_{A}f_{i}\widehat{n}_{i}\,\D A =
\iiint_{V}\partial_{i}f_{i}\,\D V.
\end{equation}
\end{example}

\subsection{Equations of motion in differential form}

\begin{node}[Mass conservation]
We see the conservation from Reynolds transport theorem, Eq~\eqref{eq:fall:lec3:mass-conservation},
\begin{equation}
0=\partial_{t}\int_{V_{c}}\rho\,\D V_{c}+\oint_{\boundary V_{c}}\rho\vec{u}\cdot\UnitNormalVec{n}\,\D A.
\end{equation}
In index notation
\begin{equation}
0=\partial_{t}\int_{V_{c}}\rho\,\D V_{c}+\oint_{\boundary V_{c}}\rho(u_{i}\widehat{n}_{i})\,\D A.
\end{equation}
Then using the divergence theorem, we have
\begin{equation}
0=\partial_{t}\int_{V_{c}}\rho\,\D V_{c}+\int_{V_{c}}\partial_{i}(\rho u_{i})\,\D V_{c}.
\end{equation}
For a \emph{stationary} control volume, this becomes
\begin{equation}
0=\int_{V_{c}}(\partial_{t}\rho + \partial_{i}(\rho u_{i}))\,\D V_{c}.
\end{equation}
But we had the control volume be \emph{completely arbitrary}. This is
a standard trick in PDEs to assert the integrand \emph{must} be
identically zero, giving us the equation:
\begin{equation}
0=\partial_{t}\rho + \partial_{i}(\rho u_{i}).
\end{equation}
This is the \define{Continuity Equation}. We could think of it as a
``pointwise'' conservation of mass. We can also think of it as using
an ``infinitesimal'' control volume.

However, this is just 1 equation in 4 unknowns (three from the
velocities $u_{i}$, and one from the density $\rho$). So this will not
have a unique solution (it could have a 3-parameter family of solutions).
\end{node}

\begin{node}[Momentum equations]
The idea is the same as before. Re-written using index notation,
\begin{equation}
\partial_{t}\iiint_{V_{c}}\rho u_{i}\,\D V_{c} + \oint_{\boundary V_{c}}\rho u_{i}u_{j}\widehat{n}_{j}\,\D A_{i}
=\iiint_{V_{c}}\gamma_{i}\,\D V_{c}+\oint_{\boundary V_{c}}f_{i}\,\D A_{c},
\end{equation}
which is really 3 equations (one for each $i=1,2,3$). Then using the
divergence theorem,
\begin{subequations}
  \begin{align}
&\partial_{t}\iiint_{V_{c}}\rho u_{i}\,\D V_{c} + \iiint_{V_{c}}\partial_{j}(\rho u_{i}u_{j})\,\D A_{i}\nonumber\\
&=\iiint_{V_{c}}\gamma_{i}\,\D V_{c} + \oint_{\boundary V_{c}}f_{i}\,\D A_{c}\\
\intertext{but assuming all surface forces are of the form $f_{i}=\sigma_{ij}n_{j}$,}
&=\iiint_{V_{c}}\gamma_{i}\,\D V_{c} + \oint_{\boundary V_{c}}\sigma_{ij}n_{j}\,\D A_{c}\\
&=\iiint_{V_{c}}(\gamma_{i} + \partial_{j}\sigma_{ij})\,\D V_{c}.
  \end{align}
\end{subequations}
We can collect terms on one side, pulling everything into one big integral:
\begin{equation}
\iiint_{V_{c}}[\partial_{t}(\rho u_{i})+\partial_{j}(\rho u_{i}u_{j})
  - \gamma_{i} - \partial_{j}\sigma_{ij}]\,\D V_{c}=0.
\end{equation}
Since the control volume was chosen arbitrarily, this means we must
have the integrand must vanish:
\begin{equation}
\partial_{t}(\rho u_{i}) + \partial_{j}(\rho u_{i}u_{j}) = \gamma_{i} + \partial_{j}\sigma_{ij}.
\end{equation}
This is the conservation of (linear) momentum equation in differential form.
\end{node}

\begin{definition}
The \define{Stress Tensor} $\sigma_{ij}$ is a symmetric tensor (has 6
independent components in $\RR^{3}$) and can be decomposed as
\begin{equation}
\sigma_{ij}=-p\delta_{i,j}+\tau_{ij}
\end{equation}
where $\tau_{ij}$ describes deviatory stresses, and $-p\delta_{i,j}$
is the pressure contribution to stresses.
\end{definition}

\begin{remark}
This means the differential form of the conservation of linear
momentum consists of 3 equations and 10 unknowns (3 from velocities
$u_{i}$, 1 from density $\rho$, and 6 from stresses $\sigma_{i,j}$).
So we have more unknowns than equations, which means this is an
under-determined system.
\end{remark}

\begin{definition}
We say an equation is in \define{Conservative Form} if we can write it as
\begin{equation}
\partial_{t}\xi + \vec{\nabla}\cdot\vec{f}(\xi)=0.
\end{equation}
If the body force has the form of $\vec{\gamma}=-\rho\vec{\nabla}V$
where $V$ is a potential, then the conservation of momentum has
conservative form.
\end{definition}